\paragraph{総消費水準目標（CLT）}
\label{sec:chapter5_a_shock_policies_clt}
CLTの式については
\ref{sec:chapter5_beta_shock_policies_clt}
を参照のこと。
\begin{enumerate}
    \item
    \label{itm:chapter5_a_shock_policies_clt_c_H_W}
        \(c_t^{H\to W}\)：
        総消費指数\(c_t^{H\to W}\)の図\ref{fig:chapter5_a_shock_graphs_c_H_W}を見ると、CLTの軌道は生産水準目標（OLT）や名目総消費水準目標（NCLT）などと並んで、ショック直後から全期間を通じて落ち込みを全政策の中で最も小さく（ほぼゼロ近傍に）抑え込んでいることがわかる。
        消費者物価水準目標（CPLT）や消費者物価インフレ目標（IT-CPI）、生産者物価水準目標（PPLT）などが初期に極端な大暴落を引き起こしているのとは対照的である。
        消費の低下は厚生を悪化させる方向に働くため、この全期間にわたる強力な消費水準の下支えが、CLTの厚生を全政策の中で最上位に押し上げる最大の要因となっている。
    \item
    \label{itm:chapter5_a_shock_policies_clt_y_H}
        \(y_t^H\)：
        生産\(y_t^H\)の図\ref{fig:chapter5_a_shock_graphs_y_H}を見ると、CLTの軌道は総消費指数\(c_t^{H\to W}\)と同様に、生産水準目標（OLT）や名目総消費水準目標（NCLT）などと並んで全期間を通じて落ち込みをほぼゼロに抑え込んでいることがわかる。
        消費者物価水準目標（CPLT）や生産者物価水準目標（PPLT）、消費者物価インフレ目標（IT-CPI）などが初期に極端な生産の大暴落を引き起こしているのとは対照的である。
        生産の低下は労働の減少を意味するため厚生にはプラスに働く。したがって生産の落ち込みを完全に防いでいるCLTは、この労働の非効用低下を通じた厚生の押し上げ効果という面では、大暴落を許容した他の政策群に比べて小さくなる。
        しかし極めて安定した生産水準を維持することで過度な変動や中盤以降の波打ちを回避しており、前述の消費の強力な下支えによる巨大なプラス効果と相まって、結果的に全政策中で最上位の厚生をもたらす要因となっている。
    \item
    \label{itm:chapter5_a_shock_policies_clt_pi_H}
        \(\pi_t^H\)：
        グロス・インフレ率\(\pi_t^H\)の図\ref{fig:chapter5_a_shock_graphs_pi_H}を見ると、CLTの定常状態からの乖離はショック直後の初期において全政策の中でも際立って大きなプラス（インフレ）方向へのスパイクを記録していることがわかる。
        生産者物価水準目標（PPLT）や生産者物価インフレ目標（IT-PPI）が全期間を通じてゼロ近傍で極めて平坦な推移を維持しているのとは対照的である。
        名目アンカーを持たず実質変数である消費のみを目標としているため、供給ショックの影響を吸収して実体経済を安定化させる過程でインフレ率の大きな上昇を許容している。この初期におけるインフレ率の大きな乖離が、次項で述べる価格分散の増大をもたらす要因となっている。
    
\end{enumerate}